\documentclass[UTF8,11pt]{ctexart}
\usepackage[margin=1in]{geometry}
\usepackage{amsmath,amssymb,amsthm,booktabs,hyperref,xcolor,enumitem,graphicx}
\emergencystretch=3em
\sloppy
\hypersetup{colorlinks=true,linkcolor=blue,urlcolor=blue,citecolor=blue}
\newtheorem{proposition}{命题}
\newtheorem{remark}{说明}
\title{用100个等圆盘覆盖单位圆盘：证书、方法与失败路线}
\author{公开计算研究档案\\\small springoxn/disk-covering-100}
\date{2026年10月}
\begin{document}
\maketitle
\begin{abstract}
本文研究用100个全等闭圆盘覆盖单位闭圆盘的最小半径。当前证书链通过六个穷尽的边界计数分支证明 $r_D(100)>0.113235$；另一个独立重放的具体构型给出 $r_D(100)\le0.1154469548855829$。目标值 $0.11324$ 尚未闭合：在当前 pair 松弛下，$B=36$ 分支仍有严格障碍。本文系统解释几何归约、有理区间验证、上下界、外部依赖以及失败路线，并把每个定量断言绑定到机器可读证据。
\end{abstract}
\section{问题与研究范围}
令 $D=\{x\in\mathbb R^2:\|x\|\le1\}$。定义
\[
r_D(100)=\inf\left\{r:\ D\subseteq\bigcup_{i=1}^{100}B(c_i,r),\quad c_i\in\mathbb R^2\right\}.
\]
这是一个连续、非凸的问题。圆盘中心不受限，覆盖诱导的 Voronoi 图或覆盖分解的组合类型会在优化过程中改变。因此，数值构型只能给出上界；下界则必须在给定半径下排除所有允许的中心配置。

本项目是面向证书的计算研究，而不是声称已经在证明助手中形式化的定理。最终判定尽量使用整数、有理数或向外取整的区间运算；浮点优化和高精度数值只用于发现候选。

\subsection{已证明与未证明}
公开发布包支持
\[
0.113235<r_D(100)\le0.1154469548855829.
\]
左侧不等式是在 $R=22647/200000$ 下的六分支排除证书；右侧不等式来自一个具体构型的重放。两者都没有确定精确最优值。特别地，附近的目标 $0.11324$ 在当前证书系统中仍然开放。

必须区分：局部盒排除不是全局证明；没有被某个障碍点阻塞不是证书；存在覆盖构型也不是最优性证明。

\section{上界：具体的100中心构型}
上界是构造性的：给出100个中心，并证明单位圆盘中每一点到某个中心的距离不超过候选半径。公开包包含两类材料。

\subsection{重放的构型}
证书 `external-disk100-bounds-67ebf1c.json` 记录了一个重放构型，给出
\[
r_D(100)\le \frac{1154469548855829}{10^{16}}.
\]
验证把单位圆盘分成边界和内部单元，对相关最远点使用精确或向外取整的算术，并绑定输入与检查器哈希。它是存在性结论，不是最优性结论。

\subsection{透明的晶格构造}
另一个证书描述了三角晶格/六角壳层构造，其候选半径 $\alpha$ 满足 $\alpha^2=1/67$。独立检查器验证保留的晶格指标、壳层计数、分离条件、正根隔离区间以及 $67t^2-1$ 的不可约性。该构造被明确作为透明的上界见证，不声称是100中心的最佳构型。

\subsection{为什么上界不是难点}
上界可以不断替换构型并重新运行覆盖检查；下界不能只优化一个有限列表，必须处理任意中心集及全部组合类型。因此主要工作集中在穷尽归约和区间证书。

\section{六分支下界架构}
固定候选半径 $R$，反设100个半径为 $R$ 的圆盘覆盖 $D$。归约按边界服务数分成
\[
B31,\ B32,\ B33,\ B34,\ B35,\ B\ge36.
\]
前五个分支利用 Euler 关系、边界所有权、三角剖分、耳、弦和内部面数 $M$，把可能的平面覆盖图化为有限模型。每个被接受的分支必须同时给出：(i) 同半径前提；(ii) 模型的穷尽覆盖；(iii) 精确或有理向外区间几何不等式；(iv) 严格为正的预算余量。

\subsection{有限组合分支}
B31--B35证书分别使用度数/曲率损失、核心耳交换、环模型、无弦模型、单弦模型，以及 $M=25$、$M=26$ 的专门归约。最终证据保存在B31、B32、B33、B34和B35 JSON证书中。这些排除结果建立在不小于当前候选半径的半径上，再由半径单调性转移到更小半径。

\subsection{高边界数分支}
对 $B\ge36$，将边界弧按循环顺序排列，并对相邻半弧参数构造局部上界。典型直接项为
\[
\Phi(a,b)=\lambda\lVert c_a-c_b\rVert^2+\frac{H(a)+H(b)}2
-\frac\mu2\bigl(\arcsin(s_a)+\arcsin(s_b)\bigr).
\]
验证器把二维参数域细分为闭有理矩形，并逐箱验证上界。循环惩罚使用精确恒等式
\[
\sum_i\frac{a_i+a_{i+1}}2=\sum_i a_i=\pi,
\]
而不是未经证明的等弧 Jensen 猜想。所得面积不等式在 $B=36$ 处有余量，并且每增加一个边界单元的增量为正，所以一旦B36 pair域闭合，尾部 $B\ge36$ 也闭合。

在 $R=22647/200000$ 下，B36重放访问1,841,089个有理箱，接受917,297个；三类标记侧支持箱数为3761、104和35。这些是证书计数，不是数值抽样。

\subsection{目标障碍}
在 $R=2831/25000=0.11324$ 下，当前 pair 松弛在 $x=y=9921033/12500000$ 附近出现精确障碍。该点的两个支持分支均为零，因此单纯细化支持区间无法修复预算缺口。这个残余是当前松弛的几何限制，而不是漏掉了网格细分。

\section{尝试过的路线：成功、失败与未决方向}
一份可信的研究记录必须包含失败路线。它们说明哪些松弛太弱，也避免重复同一障碍。

\subsection{组合与拓扑路线}
Dart map、Voronoi所有者、三角剖分和通用图编译器把几何情形转为有限组合对象。Euler计数、耳分解、非相邻耳删除和核心不等式闭合了部分低M分支。环枚举、无弦模型和单弦模型闭合了部分高M分支。主要失败点不是枚举本身，而是证明所有几何退化都被有限目录表示；因此保留模型账本和范围字段。

\subsection{容量与局部剪枝}
Hall容量树、加权见证容量、轴向多锚点盒、空间分支树、Voronoi反向距离、环带切割、径向层级、方向计数和共享顶点控制，都产生了有用的排除或候选排序。但它们单独不能闭合目标半径。非空盒仍是未解决盒，不能默认为排除。

\subsection{能量、弧与透镜路线}
角能量和谐波能量方案产生了负见证或反例，说明简单全局能量不够强。弧帽、透镜支撑、Young型支持不等式和边界弧预算构成了在0.113235闭合的B$\ge36$主路线。把粗略的 $4\lambda R^2$ 距离帽替换成精确中心距已经测试过，但该项原本就存在于直接表达式中，因此没有增加新的几何信息。共同端点透镜或真正耦合的多边联合平均仍是0.11324下B36的主要开放方向。

\subsection{半径扫描与工程化}
对0.1132301到0.11324之间的101个有理半径进行扫描，诊断了旧B36预算与精确障碍的交叉。0.113236和0.113237尝试过完整pair细分，但在停止前没有形成PASS或FAIL证书，不能把它们报告成数学失败。下一步工程改进是可续跑的有理盒队列，缓存反三角函数和支持区间，独立哈希，并按对角/非对角区域分片。

\section{软件、依赖与可复现性}
最终有理账本主要使用Python标准库，尤其是 `Fraction`、整数算术、向外取整的区间风格计算、JSON清单和SHA-256绑定。SymPy用于符号合同和因式检查；NumPy与mpmath用于发现和高精度探索；可选的python-flint/LLL工具用于部分整数关系实验。浮点结果不能单独构成最终排除证书。

项目还读取、重放并独立审计了VonEquinox的 \texttt{DiskCoveringSolve} 相关材料。它被视为外部参照和重放输入，而不是自动继承的证明。第三方源码不随本版本发布。

公开包提供 \texttt{reproduce.py}，检查核心证书文件、状态、候选半径、六分支闭合标志和SHA-256。全部线路位于 \texttt{ROUTE\_INDEX.md} 和 \texttt{route\_manifest.json}；断言与证据对应关系见 \texttt{docs/EVIDENCE\_LEDGER.md}。

PASS 具有范围：它只认证对应JSON证书中写明的域、半径、依赖和机器检查，并不表示所有几何归约都已经在证明助手中形式化。这一限制是结果的一部分，而不是脚注。

\section{证据台账与结论}
\begin{table}[h]
\centering
\small
\resizebox{\linewidth}{!}{\begin{tabular}{p{0.28\linewidth}p{0.18\linewidth}p{0.43\linewidth}}
\toprule
断言 & 状态 & 主要证据 \\
\midrule
存在100圆盘构型的上界 & 已重放 & \texttt{external-disk100-bounds-67ebf1c.json} \\
透明晶格候选 & 独立复核 & \texttt{independent\_lattice97\_upperbound\_n100.json} \\
B31--B35排除 & PASS & B31--B35完成证书 \\
$B\ge36$在$R=0.113235$下排除 & PASS & \texttt{n100-tail-candidate-224.json} \\
$R=0.113235$六分支组合 & PASS & \texttt{n100-candidate-complete-225.json} \\
$R=0.11324$旧pair松弛 & 未闭合/受阻 & support障碍和直接项诊断 \\
$R=0.113236$与$0.113237$ & 未验证 & 停止前未产生完整证书 \\
\bottomrule
\end{tabular}}
\caption{主要定量断言的状态。}
\end{table}

因此当前有证据支持的区间为
\[
\boxed{0.113235<r_D(100)\le0.1154469548855829}.
\]
左端是有证书的排除半径，右端是有证书的具体构型半径。中间差距是真正的数学工作，而不只是排版或精度问题。公开仓库的设计目标正是让这个差距及其产生路线都可检查。

\end{document}
